% LOI TUA --- Quyển XXV
\chapter*{Lời Tựa}
\addcontentsline{toc}{chapter}{Lời Tựa}
\markboth{Lời Tựa}{Lời Tựa}

\begin{truyen}{Lời Tựa}{Classroom Arguments That Stay With Us}
\chuhoa{M}{ột lớp học không chỉ có bài giảng, bài kiểm tra và điểm số. Lớp học còn có những câu hỏi ngô nghê, những màn chống chế rất nhanh trí, và những lần đối đáp khiến cả lớp vừa buồn cười vừa phải im lại để nghĩ.}
\textit{(A classroom is not only lessons, tests, and grades. It also holds naive questions, quick excuses, and exchanges that make everyone laugh first and reflect later.)}

Quyển hai mươi lăm chọn một chủ đề rất gần gũi: \textbf{lý sự cùn của học sinh}. Không phải để biến học trò thành nhân vật gây cười, mà để giữ lại những khoảnh khắc chân thật nhất trong hành trình đi học.
\textit{(Volume twenty-five chooses a very familiar theme: \textbf{student excuses and blunt arguments}. Not to turn students into jokes, but to preserve some of the most honest moments in school life.)}

Trong 10 chương và 100 màn đối đáp, bạn sẽ gặp những câu hỏi quen thuộc về việc học, điểm số, kỷ luật, tương lai, thất bại và cả những câu hỏi không có đáp án hoàn hảo. Mỗi đoạn đều ngắn, nhưng sau tiếng cười là một điều đáng giữ lại.
\textit{(Across 10 chapters and 100 exchanges, you will meet familiar questions about learning, grades, discipline, the future, failure, and even questions with no perfect answer. Each piece is short, but there is something worth keeping after the laughter.)}

Nếu bạn là học sinh, mong bạn thấy mình trong vài câu nói ở đây và hiểu rằng trưởng thành bắt đầu từ lúc biết bỏ bớt những lý lẽ dễ dãi. Nếu bạn là giáo viên, mong bạn thấy lại lớp học của mình và nhớ rằng một lời đáp đúng lúc có thể theo học trò rất lâu.
\textit{(If you are a student, may you recognize yourself here and understand that maturity begins when we let go of easy excuses. If you are a teacher, may you see your own classroom again and remember that a timely reply can stay with a learner for a long time.)}
\end{truyen}

\begin{baihoc}
Lớp học đáng nhớ không phải vì không có tiếng cười, mà vì sau tiếng cười vẫn còn lại một điều đúng.
\textit{(A memorable classroom is not one without laughter, but one where something true remains after the laughter.)}
\end{baihoc}

\begin{ghinhoanh}
\textbf{Short English to remember:}

Funny words can hide serious lessons.

A good reply can stay with a student for years.
\end{ghinhoanh}

\begin{tuhocnhanh}
1. Em từng nói câu nào trong cuốn này chưa? \textit{(Have you ever said one of these lines?)}

2. Một lời thầy cô từng nói mà em còn nhớ là gì? \textit{(What is one teacher sentence you still remember?)}

3. Sau khi đọc, em muốn bớt đi kiểu lý sự nào của mình? \textit{(Which kind of excuse do you want to reduce in yourself?)}
\end{tuhocnhanh}

\begin{flushright}
\textbf{Tôi Nguyễn Văn Sang}\\
\textit{GV THPT Nguyễn Hữu Cảnh - TP HCM}
\end{flushright}

\ngancach